\documentclass[12pt,letterpaper]{article}
\usepackage{graphicx,textcomp}
\usepackage{natbib}
\usepackage{setspace}
\usepackage{fullpage}
\usepackage{color}
\usepackage[reqno]{amsmath}
\usepackage{amsthm}
\usepackage{fancyvrb}
\usepackage{amssymb,enumerate}
\usepackage[all]{xy}
\usepackage{endnotes}
\usepackage{lscape}
\newtheorem{com}{Comment}
\usepackage{float}
\usepackage{hyperref}
\newtheorem{lem} {Lemma}
\newtheorem{prop}{Proposition}
\newtheorem{thm}{Theorem}
\newtheorem{defn}{Definition}
\newtheorem{cor}{Corollary}
\newtheorem{obs}{Observation}
\usepackage[compact]{titlesec}
\usepackage{dcolumn}
\usepackage{tikz}
\usetikzlibrary{arrows}
\usepackage{multirow}
\usepackage{xcolor}
\newcolumntype{.}{D{.}{.}{-1}}
\newcolumntype{d}[1]{D{.}{.}{#1}}
\definecolor{light-gray}{gray}{0.65}
\usepackage{url}
\usepackage{listings}
\usepackage{color}

\definecolor{codegreen}{rgb}{0,0.6,0}
\definecolor{codegray}{rgb}{0.5,0.5,0.5}
\definecolor{codepurple}{rgb}{0.58,0,0.82}
\definecolor{backcolour}{rgb}{0.95,0.95,0.92}

\lstdefinestyle{mystyle}{
	backgroundcolor=\color{backcolour},   
	commentstyle=\color{codegreen},
	keywordstyle=\color{magenta},
	numberstyle=\tiny\color{codegray},
	stringstyle=\color{codepurple},
	basicstyle=\footnotesize,
	breakatwhitespace=false,         
	breaklines=true,                 
	captionpos=b,                    
	keepspaces=true,                 
	numbers=left,                    
	numbersep=5pt,                  
	showspaces=false,                
	showstringspaces=false,
	showtabs=false,                  
	tabsize=2
}
\lstset{style=mystyle}
\newcommand{\Sref}[1]{Section~\ref{#1}}
\newtheorem{hyp}{Hypothesis}

\title{Problem Set 3}
\date{\today}
\author{Hanyu Li (25346841)}


\begin{document}
	\maketitle
	\section*{Instructions}
	\begin{itemize}
		\item Please show your work! You may lose points by simply writing in the answer. If the problem requires you to execute commands in \texttt{R}, please include the code you used to get your answers. Please also include the \texttt{.R} file that contains your code. If you are not sure if work needs to be shown for a particular problem, please ask.
		\item Your homework should be submitted electronically on GitHub in \texttt{.pdf} form.
		\item This problem set is due before 23:59 on Wednesday March 25, 2026. No late assignments will be accepted.
	\end{itemize}
	
	\vspace{.25cm}
	\section*{Question 1}
	\vspace{.25cm}
	\noindent We are interested in how governments' management of public resources impacts economic prosperity. Our data come from \href{https://www.researchgate.net/profile/Adam_Przeworski/publication/240357392_Classifying_Political_Regimes/links/0deec532194849aefa000000/Classifying-Political-Regimes.pdf}{Alvarez, Cheibub, Limongi, and Przeworski (1996)} and is labelled \texttt{gdpChange.csv} on GitHub. The dataset covers 135 countries observed between 1950 or the year of independence or the first year forwhich data on economic growth are available ("entry year"), and 1990 or the last year for which data on economic growth are available ("exit year"). The unit of analysis is a particular country during a particular year, for a total $>$ 3,500 observations. 
	
	\begin{itemize}
		\item
		Response variable: 
		\begin{itemize}
			\item \texttt{GDPWdiff}: Difference in GDP between year $t$ and $t-1$. Possible categories include: "positive", "negative", or "no change"
		\end{itemize}
		\item
		Explanatory variables: 
		\begin{itemize}
			\item
			\texttt{REG}: 1=Democracy; 0=Non-Democracy
			\item
			\texttt{OIL}: 1=if the average ratio of fuel exports to total exports in 1984-86 exceeded 50\%; 0= otherwise
		\end{itemize}
		
	\end{itemize}
	\newpage
	\noindent Please answer the following questions:
	
	\begin{enumerate}
		\item Construct and interpret an unordered multinomial logit with \texttt{GDPWdiff} as the output and "no change" as the reference category, including the estimated cutoff points and coefficients.
		\vspace{0.4cm} 
		
		\lstinputlisting[language=R, firstline=48, lastline=72]{PS3_HL.R}
		
		\medskip 
		
		Based on the research background, we have our unordered multinomial logit regression as follows:
		\[
		\ln\left(\frac{P(Y_i = \text{negative})}{P(Y_i = \text{no change})}\right) = \beta_{0, \text{negative}} + \beta_{1, \text{negative}}REG_i + \beta_{2, \text{negative}}OIL_i
		\]
		\[
		\ln\left(\frac{P(Y_i = \text{positive})}{P(Y_i = \text{no change})}\right) = \beta_{0, \text{positive}} + \beta_{1, \text{positive}}REG_i + \beta_{2, \text{positive}}OIL_i
		\]
		
		\medskip
		
		And we have the regression outcome table:
		
		\begin{table}[!htbp] \centering 
			\caption{Unordered Multinomial Logit Regression (Odds Ratios)} 
			\label{tab:unord_logit} 
			\begin{tabular}{@{\extracolsep{5pt}}lcc} 
				\\[-1.8ex]\hline 
				\hline \\[-1.8ex] 
				& \multicolumn{2}{c}{\textit{Dependent variable:}} \\ 
				\cline{2-3} 
				\\[-1.8ex] & negative & positive \\ 
				\\[-1.8ex] & (1) & (2)\\ 
				\hline \\[-1.8ex] 
				REG & 3.972$^{***}$ & 5.865$^{***}$ \\ 
				& (0.769) & (0.767) \\ 
				& & \\ 
				OIL & 119.578$^{***}$ & 97.156$^{***}$ \\ 
				& (6.885) & (6.885) \\ 
				& & \\ 
				Constant & 44.942$^{***}$ & 93.108$^{***}$ \\ 
				& (0.271) & (0.269) \\ 
				& & \\ 
				\hline \\[-1.8ex] 
				Akaike Inf. Crit. & 4,690.770 & 4,690.770 \\ 
				\hline 
				\hline \\[-1.8ex] 
				\textit{Note:}  & \multicolumn{2}{r}{$^{*}$p$<$0.1; $^{**}$p$<$0.05; $^{***}$p$<$0.01} \\ 
			\end{tabular} 
		\end{table} 
		
		\vspace{0.2cm}
		
		\begin{itemize}
			\setlength{\itemsep}{0.2cm} 
			\item Here, the two intercepts represent: for a non-democratic country (\texttt{REG = 0}) whose average ratio of fuel exports to total exports doesn’t exceed 50\% (\texttt{OIL = 0}), the odds of having positive GDP growth versus no change are expected to be $93.11$, while the odds of seeing negative GDP growth versus no change are $44.94$.
			\item Holding the fuel export ratio (\texttt{OIL}) constant, the odds of seeing negative GDP growth versus no change multiply by a factor of $3.972$ when a country is a democracy (\texttt{REG = 1}) compared to a non-democracy (\texttt{REG = 0}) and the odds of observing positive GDP growth versus no change multiply by a factor of $5.865$ when a country is a democracy compared to a non-democracy.
			\item Holding regime (\texttt{REG}) constant, the odds of seeing negative GDP growth versus no change multiply by a factor of $119.578$ when a country's average ratio of fuel exports to total exports exceeds 50\% (\texttt{OIL = 1}) compared to when this ratio is 50\% or less (\texttt{OIL = 0}). Similarly, the odds of observing positive GDP growth versus no change multiply by a factor of $97.156$ when a country's fuel export ratio surpasses 50\% compared to when it does not.
		\end{itemize}
		
		\vspace{0.4cm}
		
		To interpret the substantive impact of our predictors using predicted probabilities, we generated predictions for hypothetical cases while isolating the effect of each variable. Codes and outcomes are displayed below:
		
		\vspace{0.2cm}
		
		\lstinputlisting[language=R, firstline=79, lastline=82]{PS3_HL.R}
		
		\begin{verbatim}
			REG OIL  no change    negative   positive
			1   0   0   7.191671e-03 0.3232070  0.6696013
			2   1   0   1.378186e-03 0.2460217  0.7526001
			3   0   0   7.191671e-03 0.3232070  0.6696013
			4   0   1   6.934296e-05 0.3726529  0.6272778
		\end{verbatim} 
		
		\vspace{0.2cm}
		
		\begin{itemize}
			\setlength{\itemsep}{0.2cm}
			\item Holding the fuel export variable constant at \texttt{OIL = 0}, we compare a non-democracy (\texttt{REG = 0}) to a democracy (\texttt{REG = 1}). As the results show, when a country transitions to a democracy, the predicted probability of experiencing negative GDP growth decreases from approximately $32.3\%$ to $24.6\%$. Conversely, the predicted probability of experiencing positive GDP growth increases from $67.0\%$ to $75.3\%$.
			\item To evaluate the substantive effect of fuel exports, we hold regime type constant at non-democracy (\texttt{REG = 0}). When a country's average ratio of fuel exports to total exports exceeds 50\% (\texttt{OIL = 1}) compared to when it does not (\texttt{OIL = 0}), the predicted probability of experiencing negative GDP growth increases from approximately $32.3\%$ to $37.3\%$. Concurrently, the predicted probability of positive GDP growth decreases from $67.0\%$ to $62.7\%$.
		\end{itemize}

		\item Construct and interpret an ordered multinomial logit with \texttt{GDPWdiff} as the outcome variable, including the estimated cutoff points and coefficients.
		\vspace{0.4cm} 
		
		\lstinputlisting[language=R, firstline=84, lastline=95]{PS3_HL.R}
		\medskip 
		
		Based on the ordinal nature of our dependent variable (\texttt{negative - no change - positive}), our ordered logit regression function is specified as:
		\[
		\ln\left(\frac{P(Y_i \le j)}{1 - P(Y_i \le j)}\right) = \mu_j - (\beta_1 REG_i + \beta_2 OIL_i)
		\]
		Where $j \in \{\text{negative}, \text{no change}\}$, and $\mu_j$ represents the estimated cutoff points on the latent scale.
		
		\medskip
		
		And we have the regression outcome table:
		\begin{table}[!htbp] \centering 
			\caption{Ordered Multinomial Logit Regression(Odds Ratios)} 
			\label{} 
			\begin{tabular}{@{\extracolsep{5pt}}lc} 
				\\[-1.8ex]\hline 
				\hline \\[-1.8ex] 
				& \multicolumn{1}{c}{\textit{Dependent variable:}} \\ 
				\cline{2-2} 
				\\[-1.8ex] & GDPWdiff\_ord \\ 
				\hline \\[-1.8ex] 
				REG & 1.490$^{***}$ \\ 
				& (0.075) \\ 
				& \\ 
				OIL & 0.820$^{***}$ \\ 
				& (0.116) \\ 
				& \\ 
				\hline \\[-1.8ex] 
				$\mu_1$ (negative | no change) & -0.731 \\ 
				& \\ 
				$\mu_2$ (no change | positive) & -0.710 \\ 
				& \\
				Observations & 3,721 \\ 
				\hline 
				\hline \\[-1.8ex] 
				\textit{Note:}  & \multicolumn{1}{r}{$^{*}$p$<$0.1; $^{**}$p$<$0.05; $^{***}$p$<$0.01} \\ 
			\end{tabular} 
		\end{table} 
		
		From the regression outcomes,we can conclude that:\\
		\begin{itemize}
			\setlength{\itemsep}{0.2cm}
			\item The model estimates two latent cutoff points ($\mu_1 = -0.731$ for negative | no change, and $\mu_2 = -0.710$ for no change | positive). If a country's underlying latent economic score is below $-0.731$, the model predicts negative GDP growth; if the score falls between $-0.731$ and $-0.710$, it predicts no change; and if the score is above $ -0.710$, it predicts positive GDP growth.
			\item Holding the fuel export ratio constant, for a democratic country,the odds of having a higher degree of GDP growth are $1.490$times higher compared to a non-democratic country.
			\item Holding regime constant, the odds of moving to a higher GDP growth category multiply by a factor of $0.820$ when a country's average ratio of fuel exports to total exports exceeds 50\% compared to when this ratio is 50\% or less.
		\end{itemize}
		Whats more,we still generated fake data to predict probability and calculated probability changes:\\
		\lstinputlisting[language=R, firstline=102, lastline=111]{PS3_HL.R}
		\begin{verbatim}
			REG OIL  negative   no change   positive
			1   0   0 0.3249362 0.004555476 0.6705083
			2   1   0 0.2442235 0.003839727 0.7519368
			3   0   0 0.3249362 0.004555476 0.6705083
			4   0   1 0.3699431 0.004836151 0.6252207
		\end{verbatim} 
		\begin{table}[!htbp] \centering 
			\caption{Changes in Predicted Probabilities} 
			\label{tab:prob_changes} 
			\begin{tabular}{@{\extracolsep{5pt}}lccc} 
				\\[-1.8ex]\hline 
				\hline \\[-1.8ex] 
				& \multicolumn{3}{c}{\textit{Predicted Outcomes}} \\ 
				\cline{2-4} 
				\\[-1.8ex] \textit{Change in Variable} & negative & no change & positive \\ 
				\hline \\[-1.8ex] 
				Effect of REG (0 $\rightarrow$ 1) & $-$8.07 & $-$0.07 & +8.14 \\ 
				Effect of OIL (0 $\rightarrow$ 1) & +4.50 & +0.03 & $-$4.53 \\ 
				\hline 
				\hline \\[-1.8ex] 
			\end{tabular} 
		\end{table}
		\begin{itemize}
			\setlength{\itemsep}{0.2cm}
			\item Holding the fuel export ratio constant, transitioning from a non-democracy to a democracy decreases the probability of experiencing negative GDP growth by 8.07 percentage points and increases the probability of positive growth by 8.14 percentage points.
			\item Holding regime type constant, when a country's average ratio of fuel exports to total exports exceeds 50\%, the probability of negative GDP growth increases by 4.50 percentage points, while the probability of positive growth decreases by 4.53 percentage points.
		\end{itemize}
	\end{enumerate}
		

\section*{Question 2} 
\vspace{.25cm}

\noindent Consider the data set \texttt{MexicoMuniData.csv}, which includes municipal-level information from Mexico. The outcome of interest is the number of times the winning PAN presidential candidate in 2006 (\texttt{PAN.visits.06}) visited a district leading up to the 2009 federal elections, which is a count. Our main predictor of interest is whether the district was highly contested, or whether it was not (the PAN or their opponents have electoral security) in the previous federal elections during 2000 (\texttt{competitive.district}), which is binary (1=close/swing district, 0="safe seat"). We also include \texttt{marginality.06} (a measure of poverty) and \texttt{PAN.governor.06} (a dummy for whether the state has a PAN-affiliated governor) as additional control variables. 

\begin{enumerate}
	\item [(a)]
	Run a Poisson regression because the outcome is a count variable. Is there evidence that PAN presidential candidates visit swing districts more? Provide a test statistic and p-value.
	
	\vspace{0.4cm} 
	
	\lstinputlisting[language=R, firstline=117, lastline=128]{PS3_HL.R} 
	
	\medskip 
	
	Based on the research background and count data nature, we specified the Poisson regression model as:
	\[
	\ln(\hat{\lambda}_i) = \beta_0 + \beta_1 \text{competitive.district}_i + \beta_2 \text{marginality.06}_i + \beta_3 \text{PAN.governor.06}_i
	\]
	Where $\hat{\lambda}_i$ represents the expected mean number of visits.
	
	To test whether candidates visit swing districts more, we establish the following hypothesis for the coefficient of \texttt{competitive.district} ($\beta_1$):
	\begin{itemize}
		\item \textbf{Null Hypothesis ($H_0$):} $\beta_1 = 0$. District competitiveness has no effect on the expected number of visits.
		\item \textbf{Alternative Hypothesis ($H_a$):} $\beta_1 \neq 0$ . District competitiveness has an effect on the expected number of visits.
	\end{itemize}
	
	\medskip
	And we have the regression outcome table (with exponentiated coefficients):
	
	\vspace{0.2cm}
	\begin{verbatim}
		Coefficients:
		Estimate   Std. Error   z value   Pr(>|z|)    
		(Intercept)          -3.81023    0.22209     -17.156   <2e-16 ***
		competitive.district -0.08135    0.17069     -0.477    0.6336    
		marginality.06       -2.08014    0.11734     -17.728   <2e-16 ***
		PAN.governor.06      -0.31158    0.16673      -1.869   0.0617 .
	\end{verbatim}
	
	\begin{table}[!htbp] \centering 
		\caption{Poisson Regression (Odds Ratios)} 
		\label{tab:poi_reg} 
		\begin{tabular}{@{\extracolsep{5pt}}lc} 
			\\[-1.8ex]\hline 
			\hline \\[-1.8ex] 
			& \multicolumn{1}{c}{\textit{Dependent variable:}} \\ 
			\cline{2-2} 
			\\[-1.8ex] & PAN.visits.06 \\ 
			\hline \\[-1.8ex] 
			competitive.district & 0.922 \\ 
			& (0.171) \\ 
			& \\ 
			marginality.06 & 0.125$^{***}$ \\ 
			& (0.117) \\ 
			& \\ 
			PAN.governor.06 & 0.732$^{*}$ \\ 
			& (0.167) \\ 
			& \\ 
			Constant & 0.022$^{***}$ \\ 
			& (0.222) \\ 
			& \\ 
			\hline \\[-1.8ex] 
			Observations & 2,407 \\ 
			Log Likelihood & $-$645.606 \\ 
			Akaike Inf. Crit. & 1,299.213 \\ 
			\hline 
			\hline \\[-1.8ex] 
			\textit{Note:}  & \multicolumn{1}{r}{$^{*}$p$<$0.1; $^{**}$p$<$0.05; $^{***}$p$<$0.01} \\ 
		\end{tabular} 
	\end{table}
	
	
	\vspace{0.2cm}
	Based on the model summary, the test statistic (z-value) for \texttt{competitive.district} is $-0.477$, and the corresponding p-value is $0.6336$. Because the p-value ($0.634$) is greater than the threshold $0.05$ significance level, we don't have enough evidence to reject the null hypothesis. This means that whether a district was a swing district in the previous election isn't statistically associated with the increase of the expected number of visits from the winning PAN presidential candidate.
	
	\vspace{0.4cm}
	
	\item [(b)]
	Interpret the \texttt{marginality.06} and \texttt{PAN.governor.06} coefficients.
	
	\vspace{0.2cm}
	
	\lstinputlisting[language=R, firstline=135, lastline=136]{PS3_HL.R} 
	\begin{verbatim}
		z = 1.0668, p-value = 0.143
		alternative hypothesis: true dispersion is greater than 1
		sample estimates:
		dispersion 
		2.09834 
	\end{verbatim}
	
	\vspace{0.2cm}
	
	Before interpreting the coefficients, we conducted a dispersion test to check the equal variance assumption of the Poisson model. The test yields a p-value of $0.143$. Because $0.143 > 0.05$, we fail to reject the null hypothesis that the true dispersion is not greater than 1. Therefore, the Poisson model is correctly specified and has valid explanatory power.
	
	Therefore, we can interpret coefficients in the model as follows:
	
	\begin{itemize}
		\setlength{\itemsep}{0.2cm}
		\item \textbf{\texttt{marginality.06}:} Holding other variables constant, a one-unit increase in the marginality (poverty) index has a multiplicative effect on the expected mean of visits by $0.125$ ($e^{-2.080}$). This means that for every one-unit increase in the marginality index, we expect to see $87.5\%$ fewer visits from the winning PAN presidential candidate.
		\item \textbf{\texttt{PAN.governor.06}:} Holding other variables constant, having a PAN-affiliated governor has a multiplicative effect on the expected mean of visits by $0.732$ ($e^{-0.312}$). This implies that states with a PAN governor expect to see approximately $26.8\%$ fewer visits compared to states without a PAN governor.
	\end{itemize}
	
	\vspace{0.4cm}
	
	\item [(c)]
	Provide the estimated mean number of visits from the winning PAN presidential candidate for a hypothetical district that was competitive (\texttt{competitive.district}=1), had an average poverty level (\texttt{marginality.06} = 0), and a PAN governor (\texttt{PAN.governor.06}=1).
	
	\vspace{0.2cm}
	
	To answer this, we utilize our model to generate a fitted value for this hypothetical scenario:
	
	\vspace{0.2cm}
	
	\lstinputlisting[language=R, firstline=138, lastline=141]{PS3_HL.R} 
	
	\begin{verbatim}
		competitive.district    marginality.06   PAN.governor.06      Predicted_Visits
		       1                      0                1                 0.01494818
	\end{verbatim}
	
	\begin{itemize}
		\setlength{\itemsep}{0.2cm}
		\item Based on the prediction, for a competitive district (\texttt{competitive.district} = 1) with an average poverty level (\texttt{marginality.06} = 0) and a PAN governor (\texttt{PAN.governor.06} = 1), the estimated mean number of visits is $\hat{\lambda} \approx 0.015$. This suggests that a district with these characteristics would barely expect any visits at all from the winning candidate.
	\end{itemize}
	
\end{enumerate}

\end{document}
